\centerline{\bf \Huge Метод Шаров и перегородок}
\centerline{\bf \Huge }

\begin{flushright}
        \scriptsize{}
\end{flushright}


\problem{1}
Сколько решений имеет уравнение $x_1+x_2+\cdots+x_k=n$

a) В натуральных числах?

б) В целых неотрицательных числах?

\problem{2}
В кафе продаются пирожные четырёх сортов: эклеры, песочные, бисквитные и слоёные. Сколькими способами можно составить набор из 10 пирожных?

\problem{3}
Сколько существует пятизначных чисел, в которых цифры идут в порядке а) возрастания; б) неубывания?

\problem{4}
Сколькими способами можно разложить по 6 разным ящикам 30 одинаковых красных и 20 одинаковых синих шаров, если: а) в каждом ящике должны быть шары обоих цветов? б) нет никаких ограничений?

\problem{5}
Сколькими способами можно разделить 100 одинаковых конфет между 6 девочками, если каждой девочке должно достаться не менее 5 конфет?

\problem{6}
Сколькими способами можно выстроить в очередь 39 мальчиков и 18 девочек, чтобы никакие две девочки не стояли рядом?

\problem{7}
Сколькими способами можно расставить в ряд 5 зеленых, 30 красных и 20 синих шаров, чтобы зеленые шары не стояли рядом?

\problem{8}
а) Сколькими способами можно записать число 1000000000 в виде произведения трех натуральных чисел? Способы, отличающиеся порядком сомножителей, считаются различными.

\noindent
б) А если все множители должны быть больше 1 ?

\problem{9}
Сколько решений имеет уравнение $x+y+z+t=2020$ в нечетных натуральных числах?

\problem{10}
У Бори есть 20 жвачек, 20 сосисок и 20 гантелей. Он хочет подарить их Кириллу и Вове, каждому по 30 предметов. Сколькими способами он может это сделать?